\scanpage{58}
\mspage{-3.14-}
\margin{57}
\[
\begin{aligned}
(1,\ldots,1)\in\text{interior of }C_x
&\Longleftrightarrow
\begin{gathered}
\text{there are integers }N\ge N_0>0\\
\text{such that positive integral sums of}\\
\text{the d's generating }C_x\text{ are}\\
\text{the points:}\\
(N-N_0,N,\ldots,N),(N,N-N_0,\ldots,N),\\
\ldots,(N,N,\ldots,N-N_0)
\end{gathered}\\
&\Longleftrightarrow
P_i\text{, as in the Proposition, exist.}
\end{aligned}
\]
But on the other hand, $\{(1,1,\ldots,1)\notin\text{ interior of }C_x\}$ is
equivalent to the existence of a linear functional $\ell$ on E such that
$\ell((1,1,\ldots,1))=0$, and $\ell(d_{\alpha_0,\ldots,\alpha_n})\ge0$ for every d
generating $C_x$. Suppose
\[
\ell((z_0,\ldots,z_m))=\sum_{i=0}^m\lambda_i z_i.
\]
Without loss of generality, we may assume that $\lambda_0\le\lambda_1\le\cdots\le\lambda_m$.
The key point is:

\uline{lemma}: If for some set of $\lambda_i$ such that $\lambda_0\le\lambda_1\le\cdots\le\lambda_m$
and $\displaystyle\sum_{i=0}^m\lambda_i=0$, then
$\ell(d_{\alpha_0,\ldots,\alpha_n})\ge0$ for every $d_{\alpha_0,\ldots,\alpha_n}$
generating $C_x$, then, in fact, this happens for some $\{\lambda_i\}$ of the form:
\[
\lambda_0=\lambda_1=\cdots=\lambda_\beta=-(m-\beta)
\]
\[
\lambda_{\beta+1}=\cdots=\lambda_m=+(\beta+1).
\]
\uline{Proof}: The point is that if $\lambda_0\le\cdots\le\lambda_m$, then we can pick,
\uline{a priori}, that determinant $D_{\alpha_0,\ldots,\alpha_n}$ which is not zero at x, and
for which $\ell(d_{\alpha_0,\ldots,\alpha_n})=\displaystyle\sum_{i=0}^n\lambda_{\alpha_i}$
assumes its minimum value. Namely,
